% Input after the outer-model definitions; requires amsmath, amssymb, amsthm.
\section{外层方体径向幂权的尖锐端点}\label{sec:op-power-endpoint}
令 $\rho(x)=\|x\|_\infty$、$w_{\tau,n}(x)=\rho(x)^{\tau-n}$、
$d\mu_{\tau,n}=w_{\tau,n}\,dx$，并令 $Q(x,r)=\{y:\|y-x\|_\infty\le r\}$，
令 $H_t(E)=\{M_c\mathbf1_E>1-t\}$。本节的外层常数沿用正文定义，
其上确界只取 $0<|E|<\infty$ 及 $0<w_{\tau,n}(E)<\infty$ 的可测集。

\begin{lemma}\label{lem:op-power-tools}
对 $|F|<\infty$，有
\[|H_t(F)\setminus F|\le\frac{2^nt}{1-t}|F|.\tag{OP.2}\]
若 $d\mu_{\tau,n}=\|x\|_\infty^{\tau-n}dx$、$0<\tau<n$，则
\[\mu_{\tau,n}(B)\le K_{\tau,n}\left(|B|/2^n\right)^{\tau/n},\qquad
K_{\tau,n}=n2^n/\tau.\tag{OP.3}\]
\end{lemma}
\begin{proof}
取正文内禀上界中的 $w=1$，即得 (OP.2)。对于 (OP.3)，仅需考虑 $0<|B|<\infty$；令
$Q(0,r)$ 与 $B$ 等体积。方体内权不小于 $r^{\tau-n}$，方体外权不大于
该值；交换 $B\setminus Q(0,r)$ 与 $Q(0,r)\setminus B$ 即知
$\mu_{\tau,n}(B)\le\mu_{\tau,n}(Q(0,r))$。径向积分给出
$\mu_{\tau,n}(Q(0,r))=K_{\tau,n}r^\tau$，从而得证。
\end{proof}

\begin{theorem}[一般维数的首项]\label{thm:op-power-n-leading}
对 $0<\tau<n$，
\[C^{\rm out}_{c,w_{\tau,n}}(1-t)-1\sim t^{\tau/n}.\tag{OP.4}\]
\end{theorem}
\begin{proof}
固定 $C>2$。令 $S\subset H_t(E)\setminus E$ 为有严格见证
$r<C\|x\|_\infty$ 的短类，$T=(H_t(E)\setminus E)\setminus S$。若严格
见证恰有 $r=C\|x\|_\infty$，可微缩半径并保留严格性（$x=0$ 时等式
不可能），故 $T$ 的每个实际严格见证都满足 $r>C\|x\|_\infty$。固定半径
平均连续，故严格见证可用有理半径逼近，$S,T$ 可测。在环带
$2^k\le\|x\|_\infty<2^{k+1}$，短见证落在
$D_k=Q(0,L2^k)$，$L=2(1+C)$。于是 $S$ 与该环带之交包含于
$H_t(E\cap D_k)\setminus(E\cap D_k)$。由 (OP.2)、Tonelli 和几何级数，
\[\frac{\mu_{\tau,n}(S)}{\mu_{\tau,n}(E)}
\le\frac{2^ntL^{n-\tau}}{(1-t)(1-2^{\tau-n})}.\tag{OP.5}\]
这里对每个 $y\ne0$，环带权和至多为
$L^{n-\tau}\|y\|_\infty^{\tau-n}/(1-2^{\tau-n})$。

对长见证 $Q(x,r)$，$r>C\|x\|_\infty$，故
$Q(0,r(1-C^{-1}))\subset Q(x,r)$。见证的缺失体积小于 $t(2r)^n$，
由 (OP.3) 有
\[\mu_{\tau,n}(E)\ge K_{\tau,n}r^\tau\bigl((1-C^{-1})^\tau-t^{\tau/n}\bigr).\tag{OP.6}\]
令 $R$ 为 $T$ 中所有实际严格长见证半径的上确界（空集时 $R=0$）。
当括号为正，(OP.6) 给 $R<\infty$。$T=\varnothing$ 时长类估计为零；否则取实际见证 $(x_j,r_j)$ 使
$r_j\to R$。当 $j$ 大时，
$r_j-\|x_j\|_\infty>r_j(1-C^{-1})>R/C$，故
$Q(x_j,r_j)$ 包含 $Q(0,R/C)$。又 $T\subset Q(0,R/C)\setminus E$，所以
$|T|\le t(2R)^n$。再用 (OP.3)、(OP.6) 并令 $j\to\infty$，
\[\frac{\mu_{\tau,n}(T)}{\mu_{\tau,n}(E)}
\le\frac{t^{\tau/n}}{(1-C^{-1})^\tau-t^{\tau/n}}.\tag{OP.7}\]
取 $C=t^{-a}$、$0<a<1/n$，(OP.5) 为 $o(t^{\tau/n})$，(OP.7) 为
$t^{\tau/n}(1+o(1))$。由 $\mu(H_t(E))\le\mu(E)+\mu(S)+\mu(T)$，得到对全部 合乎定义的 输入一致的上界。

反向取 $E=Q(0,R)\setminus Q(0,h)$，其中
$h/R=t^{1/n}/(1+t^{1/n})$。对孔内 $\|x\|_\infty<h$，半径
$r=R-\|x\|_\infty$ 的方体位于 $Q(0,R)$ 并含整个内孔（小 $t$ 时
$3h<R$），其缺失比例严格小于 $t$。因此孔模零在严格光晕内，且
\[C^{\rm out}_{c,w_{\tau,n}}(1-t)-1\ge
\frac{(h/R)^\tau}{1-(h/R)^\tau}\sim t^{\tau/n}.\]
这完成证明。
\end{proof}

\begin{theorem}[一维二阶项]\label{thm:op-power-one-second}
若 $0<\tau<(5-\sqrt{17})/2$，则
\[C^{\rm out}_{c,w_{\tau,1}}(1-t)-1=t^\tau+t^{2\tau}+o(t^{2\tau}).\tag{OP.8}\]
\end{theorem}
\begin{proof}
一维短类项为
\[\frac{2t[2(1+C)]^{1-\tau}}{(1-t)(1-2^{\tau-1})}.\]
精确的一维长类估计为
\[\frac{\mu_{\tau,1}(T)}{\mu_{\tau,1}(E)}
\le\frac{2t^\tau}{A_C-2t^\tau},\qquad
A_C=(1+C^{-1})^\tau+(1-C^{-1})^\tau.\tag{OP.9}\]
事实上，对 $|x|/r<C^{-1}$，偶函数 $(1+v)^\tau+(1-v)^\tau$ 在 $v\ge0$ 上递减，故 $\mu_\tau(Q(x,r))\ge r^\tau A_C/\tau$，缺集质量至多
$2(rt)^\tau/\tau$；用实际半径上确界得到 $|T|\le2tR$，再重排即得。
取 $C=t^{-a}$。为使短类为 $o(t^{2\tau})$ 且 $C^{-2}$ 误差可忽略，取
\[\tau/2<a<(1-2\tau)/(1-\tau).\tag{OP.10}\]
此区间非空等价于 $\tau<(5-\sqrt{17})/2$，且
\[A_C=2-\tau(1-\tau)C^{-2}+O_\tau(C^{-4}),\]
所以 (OP.9) 为 $t^\tau+t^{2\tau}+o(t^{2\tau})$。
下界取 $E=[-1-t,-t]\cup[t,1+t]$。若 $0\le x<t$，取
$r=1+t-x$，则平均为 $(1-x)/(1+t-x)>1-t$；故孔模零在严格光晕内，并且
\[C^{\rm out}_{c,w_{\tau,1}}(1-t)-1\ge
\frac{t^\tau}{(1+t)^\tau-t^\tau}=t^\tau+t^{2\tau}+o(t^{2\tau}).\]
夹逼即得结论。
\end{proof}

\begin{corollary}[共同中心多集合外层模型]\label{cor:op-power-multi}
设 $a_*=\min_i a_i$。由精确外层传递
\[C^{\boldsymbol a,\rm out}_{c,w_{\tau,n}}(\lambda)=C^{\rm out}_{c,w_{\tau,n}}(1-s_{\boldsymbol a}(\lambda)),\]
有
\[C^{\boldsymbol a,\rm out}_{c,w_{\tau,n}}(1-\varepsilon)-1\sim
a_*^{-\tau/n}\varepsilon^{\tau/n}.\]
一维且 $\tau<(5-\sqrt{17})/2$ 时，左端还等于
\[a_*^{-\tau}\varepsilon^\tau+a_*^{-2\tau}\varepsilon^{2\tau}
+o(\varepsilon^{2\tau}).\]
\end{corollary}
\begin{proof}
水位方程在首段给 $s_{\boldsymbol a}(1-\varepsilon)=\varepsilon/a_*+O(\varepsilon^2)$。
代入前两定理；因 $\tau<1$，$O(\varepsilon^{\tau+1})=o(\varepsilon^{2\tau})$。
\end{proof}
